\documentclass[12pt]{article}
\usepackage[utf8]{inputenc}
\usepackage[T1]{fontenc}
\usepackage{hyperref}
\title{Timor Pidgin Language Resources for K-12 Curriculum Development}
\author{Compiled by Claude Sonnet 5 (claude-sonnet-5) for Andr\'as Kornai}
\date{}
\begin{document}
\maketitle

\section{Introduction}
Timor Pidgin (ISO 639-3 tvy) is also called Bidau Creole Portuguese. It is an extinct Portuguese-based creole once spoken in Bidau, a suburb of Dili in East Timor. Sources list it as extinct from the 1960s, when speakers shifted to standard Portuguese. The only substantial documentation found is a scholarly study based on a small corpus recorded in the early 1950s. No teaching materials of any kind turned up. This dossier is therefore short.

\section{Primary Educational Resources}
None found.

\section{Linguistic and Reference Materials}
\begin{itemize}
\item Alan N. Baxter (1990), study of the Bidau creole, in the Journal of Pidgin and Creole Languages. \url{https://www.benjamins.com/catalog/jpcl.5.1.02bax}
\item Glottolog reference entry for Baxter 1990. \url{https://glottolog.org/resource/reference/id/307143}
\item Wikipedia, Bidau Creole Portuguese. \url{https://en.wikipedia.org/wiki/Bidau_Creole_Portuguese}
\end{itemize}

\section{Religious and Community Resources}
None found.

\section{Digital and Audio Resources}
\begin{itemize}
\item Wikitongues language page for Timor Pidgin (status information only). \url{https://wikitongues.org/languages/tvy/}
\end{itemize}

\section{Curriculum Development Context}
The language is extinct and the documentation is fragmentary. A heritage curriculum could at most teach about the language and its history, using Baxter's study as the basis. Tetun-Portu and Tetun Dili are living varieties and are different from Bidau Creole Portuguese.
\end{document}